\Article{24}
Ein zweites Beispiel liefern die Dreiecke zwischen fünf Punkten der hannoverschen Triangulation: Falkenberg, Breithorn, Hauselberg, Wulfsode und Wilsede. Beobachtet wurden die Richtungen\footnote{Die Bezugsrichtungen, auf die sich die einzelnen Richtungen beziehen, gelten hier als beliebig, obwohl sie tatsächlich mit den Meridianlinien der Stationen zusammenfallen. Die Beobachtungen werden später vollständig veröffentlicht; vorerst findet sich die Figur in \textit{Astronomische Nachrichten}, Bd. I, S. 441.}:
\begin{center}
\begin{tabular}{r l l r r r}
\multicolumn{6}{c}{Bei der Station \textit{Falkenberg}}\\[5pt]
0.&Wilsede&. . . . .&\(187^0\)&\(47'\)&\(30''311\)\\
1.&Wulfsode&. . . .&225&9&39,676\\
2.&Hauselberg&. . .&266&13&56,239\\
3.&Breithorn&. . . .&274&14&43,634
\end{tabular}
\end{center}
\begin{center}
\begin{tabular}{r l l r r r}
\multicolumn{6}{c}{Bei der Station \textit{Breithorn}}\\[5pt]
4.&Falkenberg&. . .&94&33&40,755\\
5.&Hauselberg&. . .&122&51&23,054\\
6.&Wilsede&. . . . .&150&18&35,100
\end{tabular}
\end{center}
\begin{center}
\begin{tabular}{r l l r r r}
\multicolumn{6}{c}{Bei der Station \textit{Hauselberg}}\\[5pt]
7.&Falkenberg&. . .&86&29&6,872\\
8.&Wilsede&. . . . .&154&37&9,624\\
9.&Wulfsode&. . . .&189&2&56,376\\
10.&Breithorn&. . . .&302&47&37,732
\end{tabular}
\end{center}
\begin{center}
\begin{tabular}{r l l r r r}
\multicolumn{6}{c}{Bei der Station \textit{Wulfsode}}\\[5pt]
11.&Hauselberg&. . .&9&5&36,593\\
12.&Falkenberg&. . .&45&27&33,556\\
13.&Wilsede&. . . . .&118&44&13,159
\end{tabular}
\end{center}
\SourceLeaf{093}
\begin{center}
\begin{tabular}{r l l r r r}
\multicolumn{6}{c}{Bei der Station \textit{Wilsede}}\\[5pt]
14.&Falkenberg&. . .&\(7^0\)&\(51'\)&\(1''027\)\\
15.&Wulfsode&. . . .&298&29&49,519\\
16.&Breithorn&. . . .&330&3&7,392\\
17.&Hauselberg&. . .&334&25&26,746
\end{tabular}
\end{center}

Aus diesen Beobachtungen lassen sich sieben Dreiecke bilden.
\begin{center}
\begin{tabular}{l l r r r}
\multicolumn{5}{c}{Dreieck I.}\\[5pt]
Falkenberg&. . . . . .&\(8^0\)&\(0'\)&\(47''395\)\\
Breithorn&. . . . . . .&28&17&42,299\\
Hauselberg&. . . . . .&143&41&29,140
\end{tabular}
\end{center}
\begin{center}
\begin{tabular}{l l r r r}
\multicolumn{5}{c}{Dreieck II.}\\[5pt]
Falkenberg&. . . . . .&86&27&13,323\\
Breithorn&. . . . . . .&55&44&54,345\\
Wilsede&. . . . . . . .&37&47&53.635
\end{tabular}
\end{center}
\begin{center}
\begin{tabular}{l l r r r}
\multicolumn{5}{c}{Dreieck III.}\\[5pt]
Falkenberg&. . . . . .&41&4&16,563\\
Hauselberg&. . . . . .&102&33&49,504\\
Wulfsode&. . . . . . .&36&21&56,963
\end{tabular}
\end{center}
\begin{center}
\begin{tabular}{l l r r r}
\multicolumn{5}{c}{Dreieck IV.}\\[5pt]
Falkenberg&. . . . . .&78&26&25,928\\
Hauselberg&. . . . . .&68&8&2,752\\
Wilsede&. . . . . . . .&35&25&34,281
\end{tabular}
\end{center}
\begin{center}
\begin{tabular}{l l r r r}
\multicolumn{5}{c}{Dreieck V.}\\[5pt]
Falkenberg&. . . . . .&37&22&9,365\\
Wulfsode&. . . . . . .&73&16&39,603\\
Wilsede&. . . . . . . .&69&21&11,508
\end{tabular}
\end{center}
\begin{center}
\begin{tabular}{l l r r r}
\multicolumn{5}{c}{Dreieck VI.}\\[5pt]
Breithorn&. . . . . . .&27&27&12,046\\
Hauselberg&. . . . \hspace{.5em} .&148&10&28,108\\
Wilsede&. . . . . . . .&4&22&19,354
\end{tabular}
\end{center}
\SourceLeaf{094}
\begin{center}
\begin{tabular}{l l r r r}
\multicolumn{5}{c}{Dreieck VII.}\\[5pt]
Hauselberg&. . . . . .&\(34^0\)&\(25'\)&\(46''752\)\\
Wulfsode&. . . . . . .&109&38&36,566\\
Wilsede&. . . . . . . .&35&55&37,227
\end{tabular}
\end{center}

Somit liegen sieben Bedingungsgleichungen zweiter Art vor (Gleichungen erster Art entfallen offenbar). Um sie herzuleiten, sind zunächst die sphäroidischen Exzesse der sieben Dreiecke zu berechnen. Dazu muss die absolute Größe wenigstens einer Seite bekannt sein: Die Seite zwischen Wilsede und Wulfsode ist 22377,94 Meter lang. Die sphäroidischen Exzesse der Dreiecke I, II, III, IV, V, VI und VII lauten daher der Reihe nach: I\(.\,.\,.\;0''202\); II\(.\,.\,.\;2''442\); III\(.\,.\,.\;1''257\); IV\(.\,.\,.\;1''919\); V\(.\,.\,.\;1''957\); VI\(.\,.\,.\;0''321\); VII\(.\,.\,.\;1''295\).

Bezeichnen wir nun die Richtungen in der oben angegebenen Reihenfolge und mit den dort zugeordneten Indizes durch \(v^{(0)},v^{(1)},v^{(2)},v^{(3)}\) usw., so lauten die Winkel des Dreiecks I
\[
v^{(3)}-v^{(2)},\quad v^{(5)}-v^{(4)},\quad 360^0+v^{(7)}-v^{(10)}
\]
und somit lautet die erste Bedingungsgleichung
\[
-v^{(2)}+v^{(3)}-v^{(4)}+v^{(5)}+v^{(7)}-v^{(10)}+179^0\,59'\,59''798=0
\]
Entsprechend liefern die sechs übrigen Dreiecke weitere Gleichungen. Eine kurze Prüfung zeigt jedoch, dass diese sieben Gleichungen nicht unabhängig sind: Die zweite ist identisch mit der Summe der ersten, vierten und sechsten; ebenso ist die Summe der dritten und fünften identisch mit der Summe der vierten und siebten. Daher lassen wir die zweite und fünfte weg. Anstelle der verbleibenden Bedingungsgleichungen in endlicher Form schreiben wir die entsprechenden Gleichungen aus dem System (13) an und verwenden dabei (0), (1), (2) usw. für die Zeichen \(\varepsilon,\varepsilon'\) usw.:
\[
\begin{aligned}
-1''368&=-(2)+(3)-(4)+(5)+(7)-(10)\\
+1{,}773&=-(1)+(2)-(7)+(9)-(11)+(12)\\
+1{,}042&=-(0)+(2)-(7)+(8)+(14)-(17)\\
-0{,}813&=-(5)+(6)-(8)+(10)-(16)+(17)\\
-0{,}750&=-(8)+(9)-(11)+(13)-(15)+(17)
\end{aligned}
\]

Aus dem Dreieckssystem ließen sich \textit{acht} Bedingungsgleichungen dritter Art gewinnen, denn es lassen sich sowohl Dreiergruppen aus den vier Dreiecken I, II, IV und VI als auch Dreiergruppen aus den folgenden
\SourceLeaf{095}
III, IV, V und VII zu diesem Zweck miteinander kombinieren. Eine kurze Prüfung zeigt jedoch, dass \textit{zwei} genügen, eine aus jeder dieser Gruppen, da die übrigen bereits in diesen und den vorhergehenden Bedingungsgleichungen enthalten sein müssen. Die sechste Bedingungsgleichung lautet daher
\[
\begin{aligned}
&\log\sin(v^{(3)}-v^{(2)}-0''067)-\log\sin(v^{(5)}-v^{(4)}-0''067)\\
&+\log\sin(v^{(14)}-v^{(17)}-0''640)-\log\sin(v^{(2)}-v^{(0)}-0''640)\\
&+\log\sin(v^{(6)}-v^{(5)}-0''107)-\log\sin(v^{(17)}-v^{(16)}-0''107)=0
\end{aligned}
\]
und die siebte
\[
\begin{aligned}
&\log\sin(v^{(2)}-v^{(1)}-0''419)-\log\sin(v^{(12)}-v^{(11)}-0''419)\\
&+\log\sin(v^{(14)}-v^{(17)}-0''640)-\log\sin(v^{(2)}-v^{(0)}-0''640)\\
&+\log\sin(v^{(13)}-v^{(11)}-0''432)-\log\sin(v^{(17)}-v^{(15)}-0''432)=0
\end{aligned}
\]
denen die entsprechenden Gleichungen aus dem System (13) entsprechen
\[
\begin{aligned}
+25={}&+4{,}31(0)-153{,}88(2)+149{,}57(3)+39{,}11(4)-79{,}64(5)\\
&+40{,}53(6)+31{,}90(14)+275{,}39(16)-307{,}29(17)\\
-3={}&+4{,}31(0)-24{,}16(1)+19{,}85(2)+36{,}11(11)-28{,}59(12)\\
&-7{,}52(13)+31{,}90(14)+29{,}06(15)-60{,}96(17)
\end{aligned}
\]

Wenn wir nun jeder Richtung dieselbe Sicherheit zuweisen und \(p^{(0)}=p^{(1)}=p^{(2)}\ \text{etc.}=1\) setzen, und die Korrelate der sieben Bedingungsgleichungen in der hier eingehaltenen Reihenfolge mit \(A,B,C,D,E,F,G\) bezeichnen, so sind diese aus den folgenden Gleichungen zu bestimmen:
\[
\begin{aligned}
-1{,}368={}&+6A-2B-2C-2D+184{,}72F-19{,}85G\\
+1{,}773={}&-2A+6B+2C+2E-153{,}88F-20{,}69G\\
+1{,}042={}&-2A+2B+6C-2D-2E+181{,}00F+108{,}40G\\
-0{,}813={}&-2A-2C+6D+2E-462{,}51F-60{,}96G\\
-0{,}750={}&+2B-2C+2D+6E-307{,}29F-133{,}65G\\
+25={}&+184{,}72A-153{,}88B+181{,}00C-462{,}51D-307{,}29E\\
&+224868F+16694{,}1G\\
-3={}&-19{,}85A-20{,}69B+108{,}40C-60{,}96D-133{,}65E\\
&+16694{,}1F+8752{,}39G
\end{aligned}
\]

Durch Elimination erhalten wir
\SourceLeaf{096}
\[
\begin{aligned}
A&=-0{,}225\\
B&=+0{,}344\\
C&=-0{,}088\\
D&=-0{,}171\\
E&=-0{,}323\\
F&=+0{,}000215915\\
G&=-0{,}00547462
\end{aligned}
\]

Die wahrscheinlichsten Fehler ergeben sich nun aus den Formeln:
\[
\begin{aligned}
(0)&=-C+4{,}31F+4{,}31G\\
(1)&=-B-24{,}16G\\
(2)&=-A+B+C-153{,}88F+19{,}85G
\end{aligned}
\]
usw.; daraus folgen die Zahlenwerte
\[
\begin{array}{r@{\;}l@{\qquad}|@{\qquad}r@{\;}l}
(0)=&+0''065&(9)=&+0''021\\
(1)=&-0{,}212&(10)=&+0{,}054\\
(2)=&+0{,}339&(11)=&-0{,}219\\
(3)=&-0{,}193&(12)=&+0{,}501\\
(4)=&+0{,}233&(13)=&-0{,}282\\
(5)=&-0{,}071&(14)=&-0{,}256\\
(6)=&-0{,}162&(15)=&+0{,}164\\
(7)=&-0{,}481&(16)=&+0{,}230\\
(8)=&+0{,}406&(17)=&-0{,}139
\end{array}
\]

Die Summe der Quadrate dieser Fehler ergibt sich zu \(=1{,}2288\); daher beträgt der mittlere Fehler einer Richtung, soweit er sich aus den 18 beobachteten Richtungen ableiten lässt,
\[
=\BareRoot{\frac{1{,}2288}{7}}=0''4190
\]

